%% kapa: exam=1 round=19 year=2008 period=1 form=A subject=appraisal_law subject_ko=부동산관계법규 seq=40 number=80
%% source: Q-Net 감정평가사 기출문제 (cst00309) artlSeq=1001962 file=extracted/1008/1001962/2008년 제19회 감정평가사 1차시험 기출문제/1교시.hwp
%% conversion: kordoc 4.24.51 -> tools/gen1.py (auto); review=proofread-sonnet; flags=none
\begin{question}{80}
부동산등기법상 등기부에 관한 설명으로 옳지 않은 것은?
\begin{choices}
  \item 등기관은 동일한 토지에 행해진 중복등기를 발견한 경우에는 대법원규칙이 정하는 절차에 따라 중복등기되어 있는 등기용지 중 어느 하나를 남겨두고 나머지 등기용지는 쓰지 못하게 해야 한다.
  \item 등기부는 그 1용지를 등기번호란, 표제부와 갑, 을의 2구로 나누고, 각 구에는 사항란, 순위번호란을 두는데 을구는 이에 기재할 사항이 없을 때에는 이를 두지 않을 수 있다.
  \item 등기용지에는 지방법원장이 직인을 찍어야 하며, 등기부, 공동인명부와 도면은 영구히 보존하여야 한다.
  \item 누구든지 수수료를 납부하고 대법원규칙이 정하는 바에 따라 등기부의 열람 또는 그 등본이나 초본의 교부를 청구할 수 있을 뿐만 아니라, 등기부의 부속서류의 열람 및 교부를 청구할 수 있다.
  \item 대법원장은 대법원규칙이 정하는 바에 따라 등기부멸실과 관련한 멸실회복 고시에 관한 권한을 지방법원장에게 위임할 수 있다.
\end{choices}
\end{question}
\answer{4}
